\section{Degeneracy, ties and numerics}
\label{sec:robustness}

\subsection{Why the pivot loop stalls}

The normalised subproblem is degenerate by construction. In a basis the vertex
values are constant on the positive component and zero on every other tree, so a
step is blocked at length zero whenever the blocking variable is one of the many
already at zero: a vertex of the entering subtree with a negative direction
component, or a precedence slack between two vertices that are both at zero. On the
applied instances the median share of pivots with zero step is $99.3\%$, and it
reaches $99.9\%$; on the family with no arcs at all, where no slack can block,
the share is zero. Degeneracy is
therefore a property of the precedence structure, not an accident of the
implementation.

This observation drove three changes, each kept as a separate option so that
section~\ref{sec:analysis} can measure them one at a time.

\paragraph{The anti-cycling fallback must not be the operating policy.} Switching
permanently to Bland's rule after a run of degenerate pivots made every entering
rule collapse onto Bland within the first oracle, because the run of degenerate
pivots never ends. Moving the trigger far away, and releasing the fallback at the
first pivot with a positive step, restores the improving rule without giving up
termination: an episode under Bland runs on a fixed basic solution and cannot
cycle, and every episode ends with a strict increase of the ratio, so only
finitely many episodes are possible.

\paragraph{The ratio test need not scan every arc.} The direction is $+\sigma$ on
the entering subtree, $-\kappa$ on the positive component and zero elsewhere, so
an arc can block only if it crosses the boundary of one of the two sets. Because
the positive component is closed, an arc that crosses it has its tail outside,
which makes the restricted scan complete.

\paragraph{A zero step is the minimum, so the scan can stop at the first one.}
This is the shortcut the historical prototype implemented as an explicit
feasibility check during the forest visit. It is suspended while Bland's rule is
active, because that rule needs the smallest index among ties.

\subsection{Ties and numerics}

Exact ties are handled by construction: the canonical block is the union of
all maximisers, obtained either by complementarity from the dual or by
aggregating consecutive increments of equal ratio. Both settings return the same
blocks on the whole exact suite. Near-ties are separated by the relative
optimality tolerance and are deliberately not merged.

All arithmetic is \texttt{long double} with binary scaling, which is exact in
binary floating point and does not turn an integer zero-profit sum into a
spurious nonzero ratio. Certificates are re-checked in the original units, after
expanding strongly connected components, so a numerical success in the scaled
model cannot pass as a success in the user's model.
